\chapter{Data sources and preprocessing}

Different papers in this field take different approaches to obtaining their source data. One common method is to use parallel texts translated into a wide range of languages, mainly the Parallel Bible Corpus (used by \textcite{Ehret2023}). Other papers rely on annotated treebanks, such as the Universal Dependencies (used by \textcite{Coltekin2022}). This thesis aims to use as broadly multilingual data as possible and therefore relies on raw corpus frequency data.

For the purposes of this paper, three types of data are needed:
\begin{itemize}
    \item corpus data and frequencies,
    \item morphological segmentation,
    \item language-specific information.
\end{itemize}
The rest of this chapter describes how each of these data types was obtained.

\section{Corpus data}
\subsection{Raw data and aggregation}
This paper uses two main and two supplementary sources. The share of tokens contributed by each source is shown in Figure~\ref{fig:data_sources}.
\begin{itemize}

    \item Leipzig Wortschatz by \textcite{Leipzig} is the main source of corpus data. It contains corpus data for 290 languages. Most of the data come from news corpora, but Wikipedia and scraped web data are also included. The exact split is shown on the right-hand side of Figure~\ref{fig:data_sources}. The downloaded files are already frequency lists, which means that each word type appears only once together with its token count. The data were downloaded in March 2026.
    \item Glot500 by \textcite{glot500} is the second largest dataset. It contains data from 412 languages. The two datasets share $236$ languages, and $154$ languages are unique to Glot500. However, these corpora are usually smaller and of lower quality. They contain more noise in the low-frequency parts of the corpus, such as web pages, emojis, hashes, and similar artifacts.
    \item The data were further expanded using the Bible parallel corpus
    
    by \textcite{BibleCorpus}.
    \item The last source is the Universal Dependencies datasets. These datasets are annotated and well controlled, but very small.
\end{itemize}

Basic statistics of each preprocessing phase are shown in Table~\ref{tab:dataCorpusDataSizes}.

\begin{figure}[H]
    \centering
    \includegraphics[width=1\linewidth]{img/auto/data_sources.png}
    \caption{Sources of corpus data based on the file sizes. Corpus source share on right and Leipzig sources breakdown on left}
    \label{fig:data_sources}
\end{figure}

\begin{table}[H]
    \centering
    \begin{tabular}{l|r|r|r}
        Phase & Files & Total rows & Size \\
        \hline\hline
        Raw frequency data & 1 709 & 751 483 516 & 33,0 GB \\
        \hline
        Filtered data & 1 409 &  & 30,7 GB \\
        \hline
        Aggregated data & 443 & 404 637 827 & 19,7 GB \\
        \hline
        Truncated & 200 & 28 211 650 & 1,08 GB \\
        \hline
    \end{tabular}
    \caption{Size of corpus data after each preprocessing phase}
    \label{tab:dataCorpusDataSizes}
\end{table}

The data were then \textbf{aggregated} into one file per language. Token counts of word types occurring in multiple files were summed. The distribution of token counts per language is shown in Figure~\ref{fig:dist_of_total_freq}. The graph shows that half of the 443 languages contain fewer than one million tokens. The Glot500 corpus introduced many very small corpora, and 13 languages contain fewer than 100 tokens. One half of the languages lie between 162 thousand and 28 million tokens. On the other hand, 20 languages contain more than 1 billion tokens. All these statistics were calculated after the filtration described in Section~\ref{sec:filtration}.

\begin{figure}[H]
    \centering
    \includegraphics[width=1\linewidth]{img/auto/dist_of_total_freq.pdf}
    \caption{Distribution of languages with respect to their token count (log scale)}
    \label{fig:dist_of_total_freq}
\end{figure}

The data of some languages are corrupted or suboptimal for this use case. In particular:
\begin{itemize}
    \item Some languages, such as Thai and Chinese, contain whole sentences as single tokens.
    \item Some languages, such as Esperanto, have broken encoding. For Esperanto, an attempt was made to fix the data manually. Broken encoding also appears in other sources, although it does not form a substantial part of the data and can therefore be neglected. For example, in the Czech data there are 18 occurrences of the word \textit{ètyøkoalice}, which should be \textit{čtyřkoalice}, with a total frequency of 1339.
    \item The data contain further noise, described in Section~\ref{sec:data_quality}.
    \item The data may be biased by the origins of the copora making particular words as \textit{government} or \textit{president} more frequent than they would be in a more balanced corpus.
\end{itemize}

\subsection{Filtration}
\label{sec:filtration}

Several automatic measures were implemented to assess the data quality. Files that were not considered sufficiently reliable were filtered out and excluded from further processing.

\subsubsection{Corpus overlap}

Some corpora may overlap. This property was measured by the share of shared hapax legomena, computed in both directions. As explained earlier, words with very low frequencies are usually so unique that the probability of many of them occurring in multiple corpora purely by coincidence is quite small. A typical example in the Leipzig corpus is the pair of Wikipedia corpora from different years. For each file, the maximal hapax overlap with any other file was computed. A histogram of these values is shown in Figure~\ref{fig:hapax_overlap}. Most files have an overlap of about $10$ to $15\%$, which can be considered natural. Around $50\%$ there is a local maximum representing the overlapping files. The threshold should therefore be set at the local minimum between these two regions, that is, somewhere around $30\text{--}35\%$. It is reasonable to set it to exactly one third. Thus, if more than one third of the hapaxes are shared with another file, the file is considered not useful for further analysis. After removing such files, the resulting distribution is shown by the red dashed line in Figure~\ref{fig:hapax_overlap}. These files account for $2.17$~GB, which is about $6.5\%$ of the data.

\begin{figure}[H]
    \centering
    \includegraphics[width=0.8\linewidth]{img/auto/hapax_overlap.pdf}
    \caption{Distribution of the maximal hapax overlap with another file across all files. The red dashed line shows the distribution after filtration.}
    \label{fig:hapax_overlap}
\end{figure}

\subsubsection{Script unification}

Another automatic test was introduced to ensure that corpora do not mix scripts, both within a single corpus and across corpora of the same language, since some languages are commonly written in more than one script. For each file, the script distribution was calculated. Scripts were identified using the Regex library, which recognises the script from Unicode information. For efficiency reasons, only every fifth word and every third character within each sampled word, up to $10\,000$ sampled words, were considered. There were only two files in which the share of ``Unknown script'' was higher than $20\%$, namely Akha and Lahu, where the anomaly was caused by accents.

The most frequent script for a given language was taken to be the dominant one. Forty-four files were found in which the dominant script of the language covered less than $75\%$ of the characters. In the remaining files, Latin is often present as a secondary script, but the majority of words are still written in the main script of the language. There is also a gap of about $15\%$ between these 44 files and the rest, see Figure~\ref{fig:dominant_script}. These 44 files were not used in further analysis.

\begin{figure}[H]
    \centering
    \includegraphics[width=0.8\linewidth]{img/auto/dominant_script_share.pdf}
    \caption{Share of language-dominant script in the files (histogram)}
    \label{fig:dominant_script}
\end{figure}

\subsubsection{Word segmentation}

Apart from script unification, different scripts have different properties and conventions. Some scripts do not mark vowels, such as Arabic and Hebrew, while others have separate symbols for larger phonological units such as syllables, as in Korean and Devanagari. Some languages even use multiple scripts that work in different ways within a single text, as Japanese does. All these differences have to be taken into consideration in the evaluation phase. For the current phase, the important point is that some scripts do not normally separate words with spaces.

In such languages, the dataset may contain whole sentences as individual tokens. Among the first 200 languages, this applies to Chinese, Thai, Tibetan and Khmer. Chinese and Thai were segmented into words using the Python libraries \textit{jieba} and \textit{pythainlp}, respectively. For Tibetan and Khmer, no straightforward and widely used segmenters were available, and these languages were therefore excluded from the final set.

\newpage
\subsection{Truncation}
\label{sec:truncation}
The objective of this preprocessing phase is to create datasets for individual languages that are comparable regardless of their original token counts. To preserve as much information as possible, it was decided to truncate the files.

It is possible to remove words with low frequencies. Word types with very low frequencies can still be interesting, because they represent highly specific words that may not even be listed in dictionaries but are nevertheless understandable to speakers of the language. The word \textit{obviousity} (the quality of being obvious) is a good example. It is not listed in dictionaries, but its formation is clear enough that it could plausibly be used in a formal text. Even frequent spelling mistakes can be a linguistically interesting source of information.

However, in automatically collected corpora of the kind used in this analysis, most of these very low-frequency words are tokens mistakenly treated as words (such as URLs, hashes and emojis), or words from other languages. It is reasonable to assume that even correct low-frequency words are too rare for their relevance to outweigh the larger file size and the additional noise they introduce. Another important consideration is that larger datasets contain disproportionately more such items, and the goal is to make the resulting datasets as independent as possible of the original corpus size. There are several possible ways to approach the removal of rare words.
\begin{itemize}
    \item Frequency. Stop at a specific frequency (token count per word type) threshold, for example 10.
    \begin{itemize}
        \item This approach does not reflect the token count of the original file. In some languages the available data may be so limited that even important words would fall below this threshold.
    \end{itemize}
    \item Share of words. Stop at, for example, $30\%$ of all words.
    \begin{itemize}
        \item This approach does not reflect the fact that different languages need a different number of relevant words. For example, languages with rich inflection require many more word types to cover even basic communication, because the same lemma appears in multiple forms.
    \end{itemize}
    \item Top $n$ words. Stop at, for example, the 100th word.
    \begin{itemize}
        \item This approach does not reflect the original data size at all, which can also be misleading. At the same rank, a small dataset may already contain very rare words, while a large dataset may still be truncating quite frequent ones. The smallest datasets may not even contain enough words. This method also fails to reflect the structure of the language.
        \item This method was used by \textcite{Coltekin2022}, where 20~000 tokens from each language were sampled from the corresponding treebanks. The procedure was repeated 100 times to ensure stability.
    \end{itemize}
    \item Coverage. Stop when the cumulative frequency of included words reaches $95\%$ of the total frequency.
    \begin{itemize}
        \item This metric is the least dependent on the original data size. It was therefore selected, and the justification is given at the end of this chapter.
    \end{itemize}
\end{itemize}
Figure~\ref{fig:freq_rank_three} illustrates these thresholds for Spanish, Arabic and Rusyn, and shows how differently the same truncation methods behave across languages. See also the comparison of frequencies in Table~\ref{tab:threeLangCoparation}. Spanish represents one of the largest data sources, Arabic is approximately 100 times smaller, and Rusyn represents the one of the smallest usable examples. All three graphs follow Zipf's law, see Section~\ref{sec:preprocess_segmentation}.

\begin{table}[htbp]
    \centering
    \begin{tabular}{l|r|r|r}
        Language & Token count & 95\% cov. & 99\% cov. \\
        \hline\hline
        Spanish & 4 778 280 127 & 0.33 \% & 2.10 \% \\
        \hline
        Egyptian Arabic & 45 335 001 & 18.72 \% & 71.19 \% \\
        \hline
        Rusyn & 1 136 369 & 60.21 \% & 92.04 \% \\
        \hline
    \end{tabular}
    \caption{Frequencies of selected languages. Last two columns represent the percentage of the types that create 95 \% and 99 \% coverage, respectively.}
    \label{tab:threeLangCoparation}
\end{table}

\begin{figure}[p]
    \centering
    \includegraphics[width=0.75\linewidth]{img/auto/freq_rank.pdf}
    \caption{Probability of a word with respect to the ordering by frequency for three selected languages}
    \label{fig:freq_rank_three}
\end{figure}


In conclusion, the goal of this section is to obtain an approximately comparable dataset for each language. Here, ``comparable'' means that the size of the resulting dataset (the number of included words) should depend as little as possible on the original data size and as much as possible on the structure of the language itself. The graphs in Figure~\ref{fig:original_vs_truncated} show the relationship between the original size and the truncated size for all languages under the three truncation metrics. From these graphs, coverage appears to be the most suitable truncation criterion. This is also supported by Table~\ref{tab:trucation_technique_correlation}, which shows the correlation between these quantities across all languages. For $95\%$ coverage, the correlation is $0.83$, which is the lowest value.

\begin{table}[htbp]
    \centering
    \begin{tabular}{l|r|r}
        Truncation & Light & Strict \\
        \hline\hline
        Coverage $95/99$ & $0.83$ & $0.96$ \\
        \hline
        Top $5/30$ & $1.00$ & $1.00$\\
        \hline
        Frequency $100/10$ & $0.97$ & $0.98$\\
        \hline
    \end{tabular}
    \caption{Correlation between original word types and truncated token count with respect to different truncation metrics}
    \label{tab:trucation_technique_correlation}
\end{table}


\begin{figure}[p]
    \centering
    \includegraphics[width=0.7\linewidth]{img/auto/original_vs_truncated.pdf}
    \caption{Original token count vs. truncated type count with respect to different truncation methods.}
    \label{fig:original_vs_truncated}
\end{figure}

This metric has the property that the correlation is stronger for smaller datasets, where most of the words have to be included, while for larger datasets the variability is greater and depends more on the language itself. The question is whether there is a structural breakpoint on the $x$-axis (the original number of word types). Statistics for all languages and all coverage levels were calculated. The coverage levels listed in Table~\ref{tab:coverage_lvls_used} were used. The steps become smaller as the coverage approaches $100\%$, because more relevant data was expected in that region.

\begin{table}[H]
    \centering
    \begin{tabular}{l|r}
        Range & Step \\
        \hline\hline
        $50-70$ & $5$ \\
        \hline
        $72-80$ & $2$ \\
        \hline
        $81-90$ & $1$ \\
        \hline
        $91.5-100$ & $0.5$ \\
        \hline
    \end{tabular}
    \caption{Coverage levels used for analysis}
    \label{tab:coverage_lvls_used}
\end{table}

For each coverage level $c$, the procedure searches for a point such that languages with datasets larger than this point show, in total, zero correlation between the original word-type count and the truncated token count.

The model swept the candidate breakpoints $b$ (one per language) from left to right on the $x$-axis in logarithmic space and fitted a linear regression on the right segment ($\log_{10} x > b$). The right-side slope $\beta_R(b)$ starts positive for smaller corpora and ends negative, see Figure~\ref{fig:original_vs_truncated}. The first sign change, $\beta_R(b_i) \cdot \beta_R(b_{i+1}) < 0$, is detected, and the precise breakpoint with zero slope on the right side is then located by binary search between these two points.

Figure~\ref{fig:elbow} shows the location of such a point. The grey bars show the number of languages that would be included for the corresponding coverage threshold.

The relationship between the original word-type count and the truncated token count was selected, although these two metrics may seem unrelated. However, if the same metric is used on both sides, there is always a relationship between original and truncated counts that arises from the language itself. By using different metrics, it is possible to ensure that the originally collected word types do not directly influence the remaining token counts.

Finally, the elbow method was used to determine where the breakpoint starts to move rapidly towards high frequencies. The elbow method calculates the mean square error of the regression line for the points on the left $\hat y_L(c)$ and on the right $\hat y_R(c)$ of the breakpoint, and looks for the smallest deviation from those.

Let $b^*(c)$ be the optimal breakpoint for coverage level $c$. For every candidate split $c^*$, the total sum of squared errors is computed as

$$
\text{SSE}(c^*) = \sum_{c \le c^*} \bigl(\log_{10} b^*(c) - \hat{y}_L(c)\bigr)^2 + \sum_{c > c^*} \bigl(\log_{10} b^*(c) - \hat{y}_R(c)\bigr)^2
$$

The minimum of $\text{SSE}(c^*)$ is taken to be the elbow point. The smallest mean square error was found for $c=0.94$. The regression lines are plotted in green in Figure~\ref{fig:elbow}. The coverage of $94\%$ is also marked in Figure~\ref{fig:original_vs_truncated}.

\begin{figure}[htbp]
    \centering
    \includegraphics[width=0.75\linewidth]{img/auto/elbow.pdf}
    \caption{Breakpoints with zero slope on the right side for different coverage levels in blue. Red line represents the optimum found with elbow method with green slopes on both sides. The grey bars show the number of languages included for the corresponding coverage threshold.}
    \label{fig:elbow}
\end{figure}

It was therefore decided to truncate the data at $94\%$ coverage. Under this rule, languages with at least $2\ 929\ 563$ tokens are included, which is the case for $177$ languages. These $177$ languages show almost zero correlation between the original word-type count and the token count after truncation.

If truncation has to stop at a frequency at which several words share the same value, the words to be included are chosen at random. For example, if there is room for only one more word type in the dataset and $5$ word types share the same frequency of $10$, the included word is selected randomly.

\subsubsection{Published data}

The data obtained by this method will be published on Lindat. The 200 languages with the highest token counts were used. Three datasets are provided, with coverages of $90$, $94$ and $98$, in order to suit different uses of the data. The number $94$ arised from the analysis and the rest was selected symetrically from the optimum in reasonable distance. Users can voluntarily apply the same method to choose the number of languages. The breakpoints for these values are listed in Table~\ref{tab:breakpoints_for_export}.

\begin{table}[H]
    \centering
    \begin{tabular}{r|r|l}
        Coverage & Min. tokens & Recommended language count \\
        \hline\hline
        $90$ & $1\ 644\ 771$ & $199$ \\
        \hline
        $94$ & $2\ 929\ 563$ & $177$ \\
        \hline
        $98$ & $9\ 468\ 592$ & $144$ \\
        \hline
    \end{tabular}
    \caption{Recommended values for exported Lindat datasets}
    \label{tab:breakpoints_for_export}
\end{table}
\newpage
\subsection{Data quality}
\label{sec:data_quality}

As noted earlier, the datasets were collected automatically and therefore contain a substantial amount of noise. The amount of noise was measured by the following metrics:

\begin{itemize}
    \item Words written in a foreign script.
    \item Long outliers, $Q_3+3\cdot(Q_3-Q_1)$.
    \item Long outliers among the first $20$ words, $Q_3+1.5\cdot(Q_3-Q_1)$.
    \item Programming and web artifacts (URLs, emails, braces, backslash paths, multiple dots and commas).
    \item English stop words such as \textit{the, with, there, that, will, which, after, would, people, other, where, years, because, against, however}.
    \item Punctuation characters, namely \ttt{\".,?!;:/$\backslash$\@\#\$\%\^\&*$\sim$\`<>|+=\_()[]{}}.
\end{itemize}

Word types affected by the first four error types are considered corrupt. The upper graph in Figure~\ref{fig:noise} shows the proportion of corrupted types. The color indicates an approximate split between error categories, based on the values of all metrics. The bottom graph in Figure~\ref{fig:noise} shows the proportion of corrupted tokens.

For the purpose of this thesis, it was decided to keep the corrupted words, because filtering out all the noise at this scale is unrealistic. Most of the languages with a high noise rate do not use the Latin script, which suggests that languages written in Latin script are underestimated in terms of noise, since there is no easy way to detect words from other languages. The metric still serves as an indicator of how much data is, on average, invalid. Only the worst cases in terms of token share are removed. Central Khmer is removed in any case, because it is not properly tokenised into words. The other removed language is Burmese. All other languages keep the share of corrupted tokens under $5\%$.
\clearpage

\begin{figure}[p]
    \centering
    \includegraphics[width=1\linewidth]{img/auto/noise.pdf}
    \caption{Amount of corrupted word-types (upper) and tokens (lower) with error type split approximation}
    \label{fig:noise}
\end{figure}

\clearpage

\newpage
\section{Morphological segmentation data}

\subsection{Segmentation data}\label{sec:preprocess_segmentation}

For segmentation into morphemes, the model from \textcite{John2026} was used. It is a machine learning technique trained on language-specific segmented data. However, the segmentation in the source training data is not perfect, which leads to a solution of mixed quality. This task is currently under active development, and it is reasonable to expect that segmenters will improve in the near future.

The models from the \textit{morphseg} library were used. Some languages have several sources of training data with different levels of segmentation. Segmentation based on \textit{Morphynet} produces very few morphemes per word and is therefore never preferred. The same applies to the French \textit{demonette} dataset. There are three languages for which a choice between several sources had to be made. The following approach was taken. Random word types were selected from the corpus such that they differ in segmentation between the candidate sources and are valid words of the language. The segmentations were then scored on a scale from $0$, meaning completely wrong, to $1$, meaning a perfect solution, based on the author's judgement.

\begin{itemize}
    \item \textbf{English}. Based on 20 words, \textit{MorphoLex} was chosen because it scored $16$ points, while the other source scored below $11$.
    \item \textbf{German}. \textit{CELEX} is more conservative. Compared with \textit{Morphochallange}, it does not segment some obvious prefixes or German endings such as \ttt{-e}. However, when it does segment, the cut is more likely to be in the correct place. Compare the correct \ttt{knie+beug+en} with \ttt{knieb+eug+en}. Overall, the performance was very similar. \textit{CELEX} is better at identifying morphs and \textit{Morphochallange} at identifying the correct number of morphs in a word. Both were therefore kept for further phases.
    \item \textbf{Russian}. Both segmenters were also kept here. \textit{Tikhonov} performs slightly better on average, but it is more likely to make a fundamental mistake, such as not segmenting `\textrussian{поместят}' at all. \textit{KuznetsEfremDict}, on the other hand, makes more small mistakes, such as `\textrussian{на+резног+о}', where `\textrussian{ого}' should be the ending.
\end{itemize}

\subsubsection{Root identification}

The classification of morphs into roots and affixes was done using unsupervised learning. The procedure consists of two steps. In the first step, exactly one morph in each word was marked as a root. The morph with the lowest frequency in the dictionary was chosen, since roots are expected to be less frequent than affixes.

The second step handles compounds. The secondary roots are selected based on how often a particular morph is marked as a root rather than as an affix elsewhere in the data.